\documentclass[11pt]{article}
\usepackage[utf8]{inputenc}
\usepackage{amsmath,amssymb}
\usepackage{hyperref}
\title{Scalable Decision Making Neural Correlates for Large-Scale Settings}
\author{Anonymous}
\date{}
\begin{document}
\maketitle

\begin{abstract}
This paper studies Decision Making Neural Correlates. Grounded in a real, citable literature \cite{ref1,ref2,ref3,ref4,ref5,ref6,ref7,ref8},
we propose a focused method and report \textbf{illustrative} experiments.
All references below are real and verifiable (OpenAlex/arXiv); the reported
numbers are simulated to illustrate intended behaviour.
\end{abstract}

\section{Introduction}
Decision Making Neural Correlates is an active research area. This work targets a concrete gap identified from
the recent literature and contributes a method together with an evaluation protocol.

\section{Related Work (real literature)}
The following works, retrieved from public bibliographic APIs, frame our study:
\begin{itemize}
\item Craig (2008) — "How do you feel — now? The anterior insula and human awareness", Nature reviews. Neuroscience (cited 6725).
\item Bollen et al. (2011) — "Twitter mood predicts the stock market", Journal of Computational Science (cited 5163).
\item Bartra et al. (2013) — "The valuation system: A coordinate-based meta-analysis of BOLD fMRI experiments examining neural correlates of subjective value", NeuroImage (cited 2063).
\item Kable et al. (2007) — "The neural correlates of subjective value during intertemporal choice", Nature Neuroscience (cited 1979).
\item Tom et al. (2007) — "The Neural Basis of Loss Aversion in Decision-Making Under Risk", Science (cited 1767).
\item Fiske et al. (2013) — "Social Cognition: From Brains to Culture" (cited 1744).
\end{itemize}

\section{Method}
We decompose the problem across shards with a communication-efficient aggregation rule that keeps overhead sublinear in scale. We instantiate this idea for Decision Making Neural Correlates and analyse its cost/quality trade-off.

\section{Experiments (Illustrative)}
\begin{center}
\begin{tabular}{lcc}
System & Primary Metric & Efficiency \\ \hline
Strong baseline & 0.812 & 1.00$\times$ \\
Ablation & 0.828 & 0.92$\times$ \\
\textbf{Ours} & \textbf{0.851} & \textbf{0.71$\times$}
\end{tabular}
\end{center}
\emph{Metrics simulated to illustrate the intended effect; not measured.}

\section{Conclusion}
We connected a real literature to a concrete method for Decision Making Neural Correlates. Future work will
validate the illustrative results with real experiments.

\begin{thebibliography}{9}
\bibitem{ref1} A. D. Craig. How do you feel — now? The anterior insula and human awareness. \emph{Nature reviews. Neuroscience}, 2008. DOI:10.1038/nrn2555
\bibitem{ref2} Johan Bollen, Huina Mao, Xiao‐Jun Zeng. Twitter mood predicts the stock market. \emph{Journal of Computational Science}, 2011. DOI:10.1016/j.jocs.2010.12.007
\bibitem{ref3} Oscar Bartra, Joseph T. McGuire, Joseph W. Kable. The valuation system: A coordinate-based meta-analysis of BOLD fMRI experiments examining neural correlates of subjective value. \emph{NeuroImage}, 2013. DOI:10.1016/j.neuroimage.2013.02.063
\bibitem{ref4} Joseph W. Kable, Paul W. Glimcher. The neural correlates of subjective value during intertemporal choice. \emph{Nature Neuroscience}, 2007. DOI:10.1038/nn2007
\bibitem{ref5} Sabrina M. Tom, Craig R. Fox, Christopher Trepel et al.. The Neural Basis of Loss Aversion in Decision-Making Under Risk. \emph{Science}, 2007. DOI:10.1126/science.1134239
\bibitem{ref6} Susan T. Fiske, Shelley E. Taylor. Social Cognition: From Brains to Culture. \emph{}, 2013. DOI:10.4135/9781529681451
\bibitem{ref7} Jamie D. Roitman, Michael N. Shadlen. Response of Neurons in the Lateral Intraparietal Area during a Combined Visual Discrimination Reaction Time Task. \emph{Journal of Neuroscience}, 2002. DOI:10.1523/jneurosci.22-21-09475.2002
\bibitem{ref8} Wolfram Schultz. Behavioral Theories and the Neurophysiology of Reward. \emph{Annual Review of Psychology}, 2005. DOI:10.1146/annurev.psych.56.091103.070229
\end{thebibliography}
\end{document}
